\section{Рой: Swarm Kit и pioneer\_swarm}
\label{sec:swarm}

\begin{infobox}[Цели главы]
После главы вы сможете:
\begin{itemize}
    \item объяснить, как устроен рой Swarm Kit без центрального сервера: юнит, \texttt{machine\_id}, Heartbeat, почтовый ящик, подписка на телеметрию;
    \item получить список соседей и отправить сообщение конкретному дрону и всем сразу библиотекой \texttt{pioneer\_swarm};
    \item принимать сообщения без блокировки программы и придумать для них простой текстовый протокол;
    \item читать телеметрию соседа и использовать её в своём полёте: следовать за ведущим и держать место в строю;
    \item синхронизировать группу (взлёт по готовности, световое шоу) и собрать итоговый проект «поиск и спасение»;
    \item назвать ограничения: общая система координат, нет защиты от столкновений, нет гарантий доставки.
\end{itemize}
Нужно знать: полёт \texttt{pioneer\_sdk2} (глава~\ref{sec:sdk2}); главу~\ref{sec:ros2} полезно прочитать до этой (типы телеметрии те же).
\end{infobox}

\begin{missionbox}[Миссии симулятора к главе]
К главе относятся миссии 57--64 (главы трека \ui{Мини 2 · SDK2} 18--19): общение в рое (57--60) и телеметрия со строем (61--64). Миссия сама добавляет на сцену нужных соседей: Пионер Мини~2 появляются кнопками панели \ui{Рой}, у некоторых работает своя программа-бот, которую можно прочитать, выбрав дрон в панели. Программы запускаются в режиме \ui{Python · SDK2}.
\end{missionbox}

\begin{infobox}[На что опирается глава и что не проверено]
Основа главы --- документация Swarm Kit и \texttt{pioneer-swarm} на docs.geoscan.ru (Pioneer~OS 1.0.0 и новее) и десять тем в \texttt{examples/python-onboard/swarm/} (шестнадцать файлов: примеры из двух частей лежат отдельными файлами). Программы проверены на симуляторе с несколькими дронами и виртуальной сетью; \emph{на настоящих дронах авторы их не запускали}, как и на настоящей Wi-Fi-сети: в симуляторе нет задержек, потерь и перегрузки канала, доставка мгновенная. Числа сети (Heartbeat раз в секунду, исчезновение через 3~с, ящик на 100 сообщений, домены) --- допущения симулятора, справочник Геоскана их не называет; они помечены словами «в симуляторе».
\end{infobox}

\subsection{Рой без центра}
\label{sec:swarm-network}

Рой в Swarm Kit --- это несколько юнитов (бортов) в одной сети (Wi-Fi или Ethernet), которые находят друг друга, обмениваются сообщениями и делятся телеметрией \emph{напрямую}. Центрального сервера и наземной станции нет: все равноправны. У каждого юнита есть имя --- идентификатор \texttt{machine\_id} (например, \texttt{pioneer1}); по нему строятся адреса данных и почтовый ящик (рисунок~\ref{fig:swarm-mesh}).

\begin{tikzfig}
\begin{adjustbox}{max width=\linewidth}\begin{tikzpicture}[fig]
\node[forange, text width=40mm, minimum height=19mm] (a) at (0,0) {\textbf{pioneer1} \\ \footnotesize swarm-server: Heartbeat, телеметрия, ящик};
\node[forange, text width=40mm, minimum height=19mm] (b) at (10.2,0) {\textbf{pioneer2} \\ \footnotesize swarm-server: Heartbeat, телеметрия, ящик};
\node[forange, text width=40mm, minimum height=19mm] (c) at (5.1,-4.6) {\textbf{pioneer3} \\ \footnotesize swarm-server: Heartbeat, телеметрия, ящик};
\begin{scope}[on background layer]
\node[flane, fit=(a)(b)(c), inner sep=7mm, inner ysep=9mm] (net) {};
\end{scope}
\node[font=\sffamily\small\bfseries, text=Muted, anchor=north west] at ([shift={(2mm,-1.5mm)}]net.north west) {общая сеть (Wi-Fi или Ethernet), один домен};
\draw[<->, fev] (a.east) -- node[flab, above=1pt, font=\sffamily\footnotesize] {Heartbeat} (b.west);
\draw[<->, fev] ([xshift=10mm]a.south) -- node[flab, font=\sffamily\footnotesize] {Heartbeat} ([xshift=-10mm]c.north);
\draw[<->, fev] ([xshift=-10mm]b.south) -- node[flab, font=\sffamily\footnotesize] {Heartbeat} ([xshift=10mm]c.north);
\draw[fcmd, line width=1.4pt] (-1,-0.95) -- (-1,-4.6) node[flab, font=\sffamily\footnotesize, align=center, pos=0.45, left=3pt] {письмо \\ в ящик} -- (c.west);
\draw[farrow, line width=1.4pt] (11.2,-0.95) -- (11.2,-4.6) node[flab, font=\sffamily\footnotesize, align=center, pos=0.45, right=3pt] {подписка на \\ телеметрию} -- (c.east);
\end{tikzpicture}\end{adjustbox}
\captionof{figure}{Рой без центра. Каждый юнит равноправен и сам рассылает Heartbeat («я жив, вот мои данные»; штриховые стрелки); письмо адресуется в ящик конкретного \texttt{machine\_id} (оранжевая стрелка), телеметрию соседа читают по подписке (тёмная стрелка). Ни сервера, ни наземной станции нет.}\label{fig:swarm-mesh}
\end{tikzfig}


На каждом борту с Pioneer~OS работает служба \texttt{swarm-server}: она объявляет юнит в сети и отдаёт данные. Приложение на борту обращается к ней через библиотеку \texttt{pioneer\_swarm}. Юнит может быть воздушным (\texttt{AERO}), наземным (\texttt{ROVER}), водным (\texttt{MARINE}), роботом (\texttt{ROBOT}) или иным (\texttt{OTHER}); типы различаются набором объявляемых данных. У Пионера и Мини~2 в рое обычно тип воздушный: позиция, скорость, IMU, высота, батарея и события полёта.

Участники общаются тремя способами:

\begin{center}
\small
\begin{tabularx}{\linewidth}{@{}>{\raggedright\arraybackslash}p{0.23\linewidth} >{\raggedright\arraybackslash}X >{\raggedright\arraybackslash}p{0.3\linewidth}@{}}
\toprule
\textbf{Способ} & \textbf{Как устроен} & \textbf{Метод библиотеки} \\
\midrule
Обнаружение: «кто в сети?» & каждый юнит периодически рассылает Heartbeat («я жив»): свой тип и каталог данных, которые готов отдавать; из него соседи узнают, на что можно подписаться & \verb|get_unit_list()| \\
Телеметрия: «что у соседа?» & юнит подписывается на потоки соседа (позиция, батарея, \ldots); пока подписчиков нет, поток в сеть не льётся; имя потока начинается с id юнита: \fpath{pioneer1/navigation/local/position} & \texttt{get\_telemetry\_\allowbreak from\_unit()} \\
Сообщения: «написать юниту» & у юнита есть почтовый ящик: входящие строки; содержимое не регламентировано (команда, статус, текст) & \verb|send_message_to_unit()|, \verb|get_next_message()| \\
\bottomrule
\end{tabularx}
\end{center}

Если юнит пропал из сети, он исчезает из списка соседей, а подписки на него можно снять. Рой одноранговый: добавить юнит нового типа можно, не меняя остальных: из его Heartbeat понятно, что он умеет.

\begin{infobox}[Рой в симуляторе: сеть и номера]
В симуляторе \texttt{swarm-server} есть у каждого дрона сцены, работающего на \texttt{pioneer\_sdk2} (Мини~2 и Базовый с Radxa), и работает всегда --- даже без вашей программы. Дроны на \texttt{pioneer\_sdk} в рое не видны. \texttt{machine\_id} образуется из конца имени дрона («Пионер Мини~2 · 3» $\to$ \texttt{pioneer3}), а если номер занят, берётся наименьший свободный; его показывает панель \ui{Рой} под названием дрона. Heartbeat идёт раз в секунду, юнит пропадает из списка, если его Heartbeat не слышен 3~секунды (удалённый дрон, выключенный \texttt{swarm-server}). Сеть --- виртуальная: доставка мгновенная, потерь нет.
\end{infobox}

\begin{warnbox}[Кнопка «Запуск» запускает один дрон]
«Запуск» запускает программу только выбранного дрона, но программы остальных дронов продолжают работать. Чтобы запустить рой: выберите дрон, загрузите его программу, нажмите «Запуск»; затем то же на втором дроне. Поэтому примеры из двух частей лежат отдельными файлами, а порядок запуска --- \emph{сначала ожидающая сторона} (получатель, ведомый), потом отправитель или ведущий. Каждая шапка файла подсказывает, что запускать на втором дроне.
\end{warnbox}

\subsection{Библиотека pioneer\_swarm}

\begin{codebox}[Минимальная программа]
\begin{verbatim}
from pioneer_swarm import PioneerSwarmUnit, SwarmUnitType

swarm = PioneerSwarmUnit()
# ... работа с роем ...
swarm.close()
\end{verbatim}
\end{codebox}

Объект \verb|PioneerSwarmUnit| --- клиент сети в вашей программе. Создание:

\begin{codebox}[Конструктор]
\begin{verbatim}
PioneerSwarmUnit(domain_id=0, machine_id=None,
                 sync_wait_timeout_sec=None)
\end{verbatim}
\end{codebox}

\begin{center}
\small
\begin{tabularx}{\linewidth}{@{}>{\ttfamily\raggedright\arraybackslash}p{0.27\linewidth} >{\raggedright\arraybackslash}p{0.15\linewidth} X@{}}
\toprule
\textrm{\textbf{Параметр}} & \textrm{\textbf{Умолчание}} & \textrm{\textbf{Смысл}} \\
\midrule
domain\_id & 0 & домен роя: у всех юнитов должен совпадать; соседи видны только внутри домена \\
machine\_id & авто & свой идентификатор; определяется с борта сам \\
sync\_wait\_timeout\_sec & авто & сколько ждать появления id, если он ещё неизвестен \\
\bottomrule
\end{tabularx}
\end{center}

Без определённого \texttt{machine\_id} объект не создаётся (в симуляторе --- \verb|RuntimeError| с пояснением). Свойства: \verb|.id| --- свой \texttt{machine\_id}; \verb|.unit_type| --- свой тип (\verb|SwarmUnitType| или \verb|None|: значение появляется после первого собственного Heartbeat, в симуляторе --- через полсекунды после создания).

\begin{center}
\small
\begin{tabularx}{\linewidth}{@{}>{\ttfamily\raggedright\arraybackslash}p{0.43\linewidth} X@{}}
\toprule
\textrm{\textbf{Метод}} & \textrm{\textbf{Что делает}} \\
\midrule
get\_unit\_list() & список \texttt{machine\_id} соседей в сети; себя не включает \\
send\_message\_to\_unit(unit\_id, message) & отправляет строку; \verb|True| --- адресат найден и сообщение ушло, \verb|False| --- юнита в списке нет \\
get\_next\_message() & следующее входящее сообщение, \emph{без ожидания}: кортеж \verb|(sender_id, text)| или \verb|None|, если ящик пуст \\
get\_telemetry\_from\_unit(unit\_id, timeout\_sec=60.0) & ждёт телеметрию соседа и возвращает \verb|(тип, словарь потоков)|; \verb|None| --- юнит не найден \\
close() & останавливает слежение и отключается от роя \\
\bottomrule
\end{tabularx}
\end{center}

Метод \verb|close()| вызывайте, когда работа закончена, --- лучше в \verb|finally|: подписки и ящик объекта освобождаются. Опечатки в типах и значениях ловятся исключениями: неверный тип аргумента --- \verb|TypeError|, недопустимое значение (\verb|domain_id| вне 0--232, отрицательный \verb|timeout_sec|) --- \verb|ValueError|, вызов после \verb|close()| --- \verb|RuntimeError|. Список соседей в симуляторе отсортирован по номерам (\texttt{pioneer2}, \texttt{pioneer10}); у настоящей библиотеки порядок не задан, поэтому не опирайтесь на него.

\begin{infobox}[Домены и свой machine\_id]
Объект с \verb|domain_id = N| переводит юнит своего дрона в домен $N$ на время жизни объекта; соседи видны только внутри домена, письма и телеметрия в другой домен не доходят (\verb|False|, \verb|None|). Своё \verb|machine_id=...| (латинские буквы, цифры, \verb|_|) переименовывает юнит, пока объект открыт; занятое имя --- \verb|RuntimeError|. Это поведение симулятора (допущение): на настоящем борту домен и идентификатор задаёт конфигурация, сверьте её по документации Pioneer~OS.
\end{infobox}

\begin{taskbox}[Рой на панели]
Добавьте на панели \ui{Рой} ещё два Пионера Мини~2. Выпишите их \texttt{machine\_id} из панели и ответьте: кого увидит программа на первом дроне в \verb|get_unit_list()| и почему себя в списке не будет? Что изменится, если удалить второй дрон: через сколько секунд он пропадёт из списка?
\expected{у первого дрона \texttt{pioneer1}, соседи \texttt{pioneer2} и \texttt{pioneer3}; себя в списке нет по определению; удалённый дрон исчезает примерно через три секунды (три периода Heartbeat).}
\end{taskbox}
